\input{preamble}

\begin{document}

\title{K3 Surfaces}
\maketitle

\phantomsection
\label{section-phantom}
\tableofcontents

\section{K3 surfaces}
\label{section-k3-surfaces}

\begin{reference}
\cite[Definition 1.1]{huc-k3}
\end{reference}

\begin{definition}[Huybrechts]
\label{definition-k3-huybrechts}
\textit{A K3 surface over $k$} is a complete
non-singular variety $X$ of dimension two
such that
$$
\Omega^2_{X/k}\simeq\mathcal{O}_X.
$$
and
$$
H^1(X,\mathcal{O}_X)=0.
$$
\end{definition}

\noindent
By definition, a variety over a field $k$ is
complete if the morphism
$X\to\operatorname{Spec}(k)$ is proper.
The variety $X$ over $k$ is non-singular if
the cotangent sheaf $\Omega_{X/k}$ is locally
free of rank $\dim(X)$. This is equivalent
to
$X_{\overline{k}}=X\times_k\overline{k}$
being regular; see, for example, [375,
Proposition 6.2.2].

\begin{definition}[Misha]
\label{definition-k3-misha}
\textit{A K3 surface} is a compact complex
surface $M$
with
$b_1(M)=0$ and
$c_1(M,\mathbb{Z})=0$.
\end{definition}

\begin{remark}
This implies that the canonical
bundle is trivial.
\end{remark}

\begin{proposition}
The Hodge diamond of a K3 surface
is
\[
\begin{matrix}
&&1&&\\
&0&&0&\\
1&&20&&1\\
&0&&0&\\
&&1&&
\end{matrix}
\]
so in particular
$H^{2,0}(M)$ is one-dimensional.
\end{proposition}

\begin{remark}
The middle cohomology is the main
cohomological invariant of a K3
surface.
The intersection form lives on
$H^2(M,\mathbb{Z})$.
\end{remark}

\medskip\noindent
In this exercise we shall find
an example of a non-algebraic K3 surface.
Recall that the algebraic dimension $a(X)$
of a variety $X$ is the transcendence
degree of its field of meromorphic
functions over $\mathbb{C}$.

\begin{exercise}
\label{exercise-k3-algebraic-dimension}
Prove that there exists a K3 surface
of algebraic dimension 1.
\end{exercise}

\begin{proof}
The main idea is that a surface
is algebraic if and only if $a(X) = 2$.
We need to construct a surjection
$f:X \to \mathbb{P}^1$
for some nonprojective $X$.

Pulling back
$\mathcal{M}(\mathbb{P}^1)=\mathbb{C}(t)$
to $X$ we obtain an injection
$f^*\mathbb{C}(t)\hookrightarrow\mathcal{M}(X)$,
and so $a(X) \geq 1$.

Showing that $X$ is nonprojective
thus gives $a(X) = 1$.

My idea so far is to use a Kummer surface
$K$. For starters, since $K$ is defined
as the normalization of an abelian
surface $A$ quotiented by an involution
$\iota$, we can deal only with $A/\iota$
since normalization is a dominant map.

Then we argue that quotient by $\iota$
also makes no difference in
$\mathcal{M}(A)$.

But I'm not sure how to create the map
$K \to \mathbb{P}^1$, or
$A \to \mathbb{P}^1$.

Also not sure how to prove $K$ is not
projective.

\medskip\noindent
{\bf Day 2.} Today I did the following:
let $E = \mathbb{C}/\mathbb{Z}^2$.
Regular functions on $E$ are double
periodic functions, i.e. such that
$f(z) = f(z+n e_1+ m e_2)
=f(z+n+im)$ for all $n,m \in \mathbb{Z}$.
They are all constant by Liouville.
But there can be doubly periodic
{\it meromorphic} functions: having poles
makes them be unbounded.

In fact, there exist nontrivial meromorphic
functions such as the Weierstrass
something function.
Being nontrivial, it has positive
transcendence degree, for no nonconstant
function can be constructed from 
complex numbers.

Now consider $E \times E \to E$.
I aim to show that $E \times E$ is not
projective. By Huybrechts HK paper,
it's enough to show that the intersection
form is negative semidefinite, i.e. that
$L^2 \leq 0$ for all $L \in
\operatorname{Pic}_{E \times E}$,
which is a criterion that works for
hyperkähler manifolds such as the K3
surface $E \times E$.

I still need to use (and prove)
the criterion that
projective iff $a(X)=2$.
So probably: if projective,
restrict $\mathcal{M}(\mathbb{P}^K)=
\mathbb{C}(t_1,…)$ to $X$,
this gives $a(X) \geq  2$.
But $a(X)$ cannot be greater than 2
if $X$ is a surface.
So: not projective, not 2; surface, not
$\geq 2$. Elliptic: $a(X) \geq 1$.

\medskip\noindent
{\bf Day 3.} Maybe I should just
try to prove that $a(E\times E)<2$
somehow…
\end{proof}

\section{Curves on K3 surfaces}
\label{section-curves-on-k3}

\noindent
My progress on this
exercise still has three black boxes:
\begin{enumerate}
\item Construction of normalization
from gluing affine pieces.
\item Finite scheme morphism preserves Euler
characteristic.
 
I know that the normalization map 
$\nu:\widetilde C\to C$
is finite by Noether normalization
(and gluing, which is my current work).
The roadmap is:
Finite morphism $\implies$
affine morphism $\implies$
preserves cohomology, i.e.
for any sheaf $\mathcal{F}$ on
$\widetilde C$ we have
$$
H^i(\widetilde C, \mathcal{F})
\cong H^i(C,\nu_*\mathcal{F}).
$$
That uses Leray Spectral Sequence,
apparently.
But it obviously gives
$\chi(\mathcal{O}_{\widetilde{C}})
=\chi (\nu_*\mathcal{O}_{\widetilde{C}})$
which is what we really need.
\item If the local ring is regular
i.e. $\mathfrak{m}/\mathfrak{m}^2
=\dim \mathcal{O}_p$
then the variety is smooth
in the sense that it has 
holomorphic coordinates near $p$.
If I presented this on Friday I would have to
prove that.
\end{enumerate}

\begin{exercise}
\label{exercise-singular-curve-on-k3}
Let $C$ be a singular, irreducible complex
curve on a K3 surface $X$. Prove that
$(C.C)\geq 0$.
\end{exercise}

\begin{proof}
By the Genus Formula
(Lemma
\ref{complex-geometry-lemma-genus-formula})
we have
$$
2p_a(C)-2 = C \cdot (C + K_X).
$$
Since $X$ is a K3 surface, the canonical
divisor is linearly equivalent to zero.
Hence
$$
C^2 = 2p_a(C)-2.
$$
Thus it suffices to prove that
$p_a(C) \geq 1$.

\medskip\noindent
Let
$$
\nu : \widetilde C \longrightarrow C
$$
be the normalization of $C$. We recall its
construction. Choose an affine open covering
$U_i = \Spec(A_i)$ of $C$. Since $C$ is
irreducible, all the rings $A_i$ have the
same function field $K(C)$. Let
$\overline A_i$ be the integral closure of
$A_i$ in $K(C)$. Then the affine schemes
$$
\Spec(\overline A_i)
$$
glue to a normal curve $\widetilde C$,
and the inclusions
$$
A_i \subset \overline A_i
$$
glue to a finite birational morphism
$\nu : \widetilde C \to C$
(finiteness follows
from Noether normalization lemma,
i.e. Commutative Algebra
Lemma \ref{commutative-algebra-lemma-noether-normalization},
and birationality is basically by
construction of the normalization
as the spectrum of the integral
closure, i.e. the fraction fields
are the same and the induced
map on the original ring is the identity;
check the definition of birational
morphism, Commutative Algebra Definition
\ref{commutative-algebra-definition-birational-morphism}.).

We have an exact sequence of sheaves
$$
0 \to \mathcal O_C
  \to \nu_*\mathcal O_{\widetilde C}
  \to \delta \to 0,
$$
where
$$
\delta =
\nu_*\mathcal O_{\widetilde C}/\mathcal O_C.
$$
Taking Euler characteristics gives
$$
\chi(\mathcal O_C)
=
\chi(\nu_*\mathcal O_{\widetilde C})
-
\chi(\delta).
$$
Since $\nu$ is finite (this is
a black box for now)
$$
\chi(\nu_*\mathcal O_{\widetilde C})
=
\chi(\mathcal O_{\widetilde C}).
$$

Now observe that, at smooth
points $p \in C$, the stalks
$\mathcal{O}_{C,p}$ are
UFDs by Complex Analysis
Lemma \ref{complex-analysis-lemma-stalk-is-ufd}
--- due to Weierstrass Preparation
Theorem, in turn due to the Gauss lemma.
By Commutative Algebra Lemma
\ref{commutative-algebra-proposition-ufd-integrally-closed},
UFDs are integrally closed,
which means that the stalks
of $\tilde{C}$ and $C$ are isomorphic.
Thus, the stalks of $\delta$
can only be non trivial at non smooth
points of $C$. There can only
be finitely many non smooth points
in a projective variety,
and therefore $\delta$
is a skyscraper sheaf with finite support.
Thus
$$
\chi(\delta)=h^0(\delta).
$$
Therefore
$$
\begin{aligned}
p_a(C)
&= 1-\chi(\mathcal O_C) \\
&= 1-\chi(\mathcal O_{\widetilde C})
   + h^0(\delta) \\
&= g(\widetilde C)+h^0(\delta),
\end{aligned}
$$
where in the last equality
we we can replace
$1-\chi(\mathcal{O}_{\widetilde{C}})$
with the geometric genus
$g(\widetilde{C}):=
h^0(\omega_{\widetilde{C}})$
by applying Serre duality,
thereby using that $\widetilde{C}$
is smooth since it is normal.
To prove that normal implies
smooth we use
Commutative Algebra
Proposition
\ref{commutative-algebra-proposition-integrally-closed-dvr}
to get from integrally closed
that the local ring is regular.
Then we need Hartshorne
to get differential-geometric smoothness
from regular local ring.

Now $g(\widetilde C)\geq 0$. Since $C$ is
singular, the normalization is not an
isomorphism, so $\delta\neq 0$. Since
$\delta$ is a nonzero skyscraper sheaf, it
has a nonzero global section. Hence
$$
h^0(\delta)\geq 1.
$$
Thus
$$
p_a(C)=g(\widetilde C)+h^0(\delta)\geq 1.
$$
Consequently
$$
C^2=2p_a(C)-2\geq 0.
$$
\end{proof}



\section{Period domain}
\label{section-period-domain}

\begin{definition}
\label{definition-k3-lattice}
The {\it K3 lattice} is the even
unimodular lattice
$$
\Lambda_{K3}
=
U^{\oplus 3}\oplus E_8(-1)^{\oplus 2}
$$
of signature $(3,19)$, where $U$ is the
hyperbolic plane with intersection matrix
$$
\begin{pmatrix}
0&1\\
1&0
\end{pmatrix},
$$
and $E_8(-1)$ is the negative-definite
$E_8$ lattice.
For every K3 surface $X$, the lattice
$H^2(X,\mathbb{Z})$ equipped with its
intersection pairing is isomorphic to
$\Lambda_{K3}$.
\end{definition}

\begin{definition}
\label{definition-period-k3-surface}
Let $X$ be a K3 surface.
The {\it period} of $X$ is the
weight-two integral Hodge structure
$$
H^2(X,\mathbb{C})
=
H^{2,0}(X)\oplus H^{1,1}(X)
\oplus H^{0,2}(X)
$$
on $H^2(X,\mathbb{Z})$, together with
the intersection pairing
$$
(\alpha,\beta)=\int_X\alpha\wedge\beta.
$$
Equivalently, after choosing a {\it
marking}, i.e. a lattice isomorphism
$\phi:H^2(X,\mathbb{Z})\to\Lambda_{K3}$,
the period is the line
$$
\phi_{\mathbb{C}}(H^{2,0}(X))
\subset\Lambda_{K3}\otimes\mathbb{C}.
$$
\end{definition}

\begin{definition}
\label{definition-period-domain-k3-surfaces}
Let $\Lambda_{K3}$ be the K3 lattice.
The {\it period domain} of marked K3
surfaces is
$$
\Omega_{\Lambda_{K3}}
=
\left\{
[\sigma]\in
\mathbb{P}(\Lambda_{K3}\otimes\mathbb{C})
\ \middle|\
(\sigma,\sigma)=0,\quad
(\sigma,\overline{\sigma})>0
\right\}.
$$
Thus $\Omega_{\Lambda_{K3}}$ parametrizes
the weight-two Hodge structures on
$\Lambda_{K3}$ which are compatible with
its intersection pairing and have
one-dimensional $(2,0)$-part.
\end{definition}

\begin{definition}[Misha]
Let $V=H^2(M,\mathbb{R})$ with the
intersection form.
The \textit{period domain} is
\[
\mathbb{P}\mathrm{er}=
\{[v]\in \mathbb{P}H^2(M,\mathbb{C})
\mid (v,v)=0,\ (v,\bar{v})>0\}.
\]
\end{definition}

\begin{remark}
This is the open subset of the
quadric cut out by $(v,v)=0$.
The positivity condition chooses
the correct connected component.
\end{remark}

\begin{proposition}[Misha]
The period point of a K3 surface
is the line $H^{2,0}(M)$.
\end{proposition}

\begin{remark}
So the period map sends a complex
structure to its Hodge line in
$H^2(M,\mathbb{C})$.
The old notes identify this
space with a Grassmannian of
positive oriented planes.
\end{remark}

\section{Local Torelli}
\label{section-local-torelli}

\begin{definition}[Misha]
\label{definition-period-map-misha}
Let $\mathrm{Comp}$ be the space of complex
structures $I$ on $M$ for which $(M,I)$ is
a K3 surface, and let $\mathrm{Diff}_0(M)$
be the connected component of the identity
in the diffeomorphism group of $M$.
Set
$$
\mathrm{Teich}
=
\mathrm{Comp}/\mathrm{Diff}_0(M).
$$
The \textit{period map} is
$$
\begin{aligned}
\mathrm{Per}:\mathrm{Teich}
&\longrightarrow \mathbb{P}\mathrm{er},\\
I&\longmapsto H^{2,0}(M,I).
\end{aligned}
$$
\end{definition}

\begin{proposition}
The period map for K3 surfaces is
a local diffeomorphism.
\end{proposition}

\begin{remark}
This is the deformation-theoretic
statement that connects complex
structure with the Hodge subspace
$H^{2,0}(M)$.
\end{remark}

\begin{proposition}
On a K3 surface, a holomorphic
symplectic form determines the
complex structure locally.
\end{proposition}

\section{Lefschetz}
\label{section-lefschetz}

\begin{remark}
Morse theory enters the proof of
Lefschetz hyperplane section via
a strictly plurisubharmonic
function on the affine complement.
\end{remark}

\begin{proposition}
If $Z\subset \mathbb{C}P^n$ is
smooth and $H$ is a hyperplane,
then the topology of
$Z\setminus H$ is controlled by
the critical points of an
appropriate Morse function.
\end{proposition}

\begin{remark}
The outline is:
\begin{enumerate}
\item Work on the affine
complement $\mathbb{C}P^n\setminus H$.
Take the standard potential
$f_1=\sum |z_i|^2$ on
$\mathbb{C}^n$.
\item Its complex Hessian
$dd^cf_1$ is positive.
So $f_1$ is strictly
plurisubharmonic.
\item Perturb $f_1$ slightly.
Strict plurisubharmonicity is
open in the $C^2$ topology and
Morse functions are dense.
Hence choose a function that is
both Morse and strictly
plurisubharmonic on
$Z\setminus (Z\cap H)$.
\item Use the gradient flow of
that function.
Each critical point gives a
stable manifold.
The stable manifolds are the
pieces attached to the lower
level set.
\item The stable manifolds have
real dimension equal to the
Morse index, and for a
strictly plurisubharmonic
function the index is at most
$\dim_{\mathbb{C}} Z$.
\item Therefore the open
complement is built from the
hyperplane section by attaching
handles of controlled index.
This gives the desired
connectivity and cohomological
vanishing range.
\end{enumerate}
\end{remark}

\begin{remark}
This is the point of the old
notes: one identifies
the hyperplane section as a
deformation retract of the
complement of the stable
manifolds.
That turns the Lefschetz
statement into a Morse
computation.
\end{remark}

\begin{thebibliography}{9}
\bibitem{huyk3}
D. Huybrechts,
\textit{Lectures on K3 Surfaces}.
\end{thebibliography}

\end{document}
